\subsection{Limits of measurable functions}

THEOREM 2 (Egoroff). --- Let X be a locally compact space, \(\mu\) a measure on X, A a countable set, F a filter on A having a countable base, and \((f_{\alpha})_{\alpha\in\mathbb{A}}\) a family of measurable mappings of X into a metrizable space F. Assume that \(\lim_{\mathfrak{F}} f_{\alpha}(x) = f(x)\) exists in the complement of a locally negligible subset N of X. Under these conditions,

\(1^{\circ}\) the function \(f\) (extended to \(\mathbf{N}\) in any manner whatsoever) is measurable;

\(2^{\circ}\) for every compact subset K of X and every \(\varepsilon > 0\) , there exists a compact set \(\mathrm{K}_{1} \subset \mathrm{K}\) such that \(|\mu| (\mathrm{K} - \mathrm{K}_{1}) \leqslant \varepsilon\) and such that the restrictions of the \(f_{\alpha}\) to \(\mathrm{K}_{1}\) are continuous and converge uniformly to \(f\) on \(\mathrm{K}_{1}\) .

The first assertion obviously follows from the second, which we are going to prove. There exists a compact set \(K_{0} \subset K\) such that \(|\mu|(\mathrm{K} - \mathrm{K}_{0}) \leqslant \varepsilon/2\) and such that the restrictions to \(K_{0}\) of all the functions \(f_{\alpha}\) are continuous (No.~1, Prop. 2). Let \((\mathrm{A}_{n})\) be a decreasing countable base for the filter F; let d be a metric on F compatible with the topology. For every pair of integers \(n > 0\) , \(r > 0\) , let \(\mathbf{B}_{n,r}\) be the set of points \(x \in \mathrm{K}_0\) such that, for at least one pair of indices \(\alpha, \beta\) belonging to \(\mathbf{A}_n\) , \(d(f_{\alpha}(x), f_{\beta}(x)) \geqslant 1 / r\) ; for fixed \(\alpha\) and \(\beta\) , the set of \(x \in \mathrm{K}_0\) such that \(d(f_{\alpha}(x), f_{\beta}(x)) \geqslant 1 / r\) is closed in \(\mathbf{K}_0\) , hence is compact; consequently, \(\mathbf{B}_{n,r}\) is a countable union of compact sets contained in \(\mathbf{K}_0\) , hence is integrable (§4, No.~5, Props. 6 and 8). If \(r\) is fixed, the intersection of the decreasing sequence of sets \(\mathbf{B}_{n,r}\) ( \(n = 1,2,\ldots\) ) has measure zero, since \(f_{\alpha}(x)\) tends to \(f(x)\) almost everywhere in \(\mathbf{K}_0\) with respect to the filter \(\mathfrak{F}\) ; thus \(\lim_{n \to \infty} |\mu|(\mathbf{B}_{n,r}) = 0\) (§4, No.~5, Cor. of Prop. 7), consequently there exists an integer \(n_r\) such that \(|\mu|(\mathbf{B}_{n_r,r}) \leqslant \varepsilon / 2^{r+2}\) . Let \(B\) be the union (for \(r = 1,2,\ldots\) ) of the sets \(\mathbf{B}_{n_r,r}\) ; \(B\) is integrable and

\[
| \mu | (\mathrm{B}) \leqslant \sum_ {r = 1} ^ {\infty} | \mu | (\mathrm{B} _ {n _ {r}, r}) \leqslant \varepsilon / 4
\]

(§4, No.~5, Cor. of Prop. 8). Let C be the complement of B in \(K_{0}\) ; by construction, \(f_{\alpha}(x)\) converges uniformly to \(f(x)\) in C with respect to the filter F, and since the restrictions of the \(f_{\alpha}\) to C are continuous, so is the restriction of f to C. It then suffices to take a compact set \(K_{1} \subset C\) such that \(|\mu|(\mathrm{C} - \mathrm{K}_{1}) \leqslant \varepsilon/4\) in order to satisfy the conditions of the statement, since \(|\mu|(\mathrm{K} - \mathrm{K}_{1}) = |\mu|(\mathrm{K} - \mathrm{K}_{0}) + |\mu|(\mathrm{B}) + |\mu|(\mathrm{C} - \mathrm{K}_{1}) \leqslant \varepsilon\) .

The conclusions of Th. 2 do not necessarily hold if F is not metrizable (Exer. 1). If F is metrizable, and the set A is not countable but the filter \(\mathfrak{F}\) has a countable base, then the first conclusion of Th. 2 is again valid; for, if \((\mathrm{A}_n)\) is a countable base for \(\mathfrak{F}\) and \(\alpha_{n}\) is an element of \(\mathrm{A}_n\) , then \(f\) is the limit of the sequence \((f_{\alpha_n})\) locally almost everywhere, hence is measurable; however, the second conclusion of Th. 2 is no longer necessarily valid (cf.~Exer. 4).

COROLLARY 1. --- Let \((f_n)\) be a sequence of numerical functions (finite or not). If the \(f_n\) are measurable, then the functions \(\sup_{n} f_n\) , \(\inf_{n} f_n\) , \(\limsup_{n \to \infty} f_n\) and \(\liminf_{n \to \infty} f_n\) are measurable.

For, the extended real line \(\overline{\mathbf{R}}\) , being homeomorphic to a compact interval of \(\mathbf{R}\) , is metrizable. The function \(\sup_{n} f_n\) is the pointwise limit of the increasing sequence of functions \(g_n = \sup(f_1, f_2, \ldots, f_n)\) , which are measurable (No.~3, Cor. 1 of Th. 1); similarly, \(\limsup_{n \to \infty} f_n\) is the pointwise limit of the decreasing sequence of functions \(h_n = \sup_{p \geqslant 0} f_{n+p}\) , each of which is measurable by the foregoing. Finally, since \(\inf_{n} f_n = -\sup_{n} (-f_n)\) and \(\liminf_{n \to \infty} f_n = -\limsup_{n \to \infty} (-f_n)\) , these functions are measurable.

COROLLARY 2. --- The measurable sets in X form a tribe (GT, IX, §6, No.~3).

For, if M and N are measurable then the functions

\[
\varphi_ {\mathrm{M} \cup \mathrm{N}} = \varphi_ {\mathrm{M}} + \varphi_ {\mathrm{N}} - \varphi_ {\mathrm{M}} \varphi_ {\mathrm{N}} \quad \text {and} \quad \varphi_ {\mathrm{M} \cap \mathbf {C N}} = \varphi_ {\mathrm{M}} - \varphi_ {\mathrm{M}} \varphi_ {\mathrm{N}}
\]

are measurable by No.~3, Cors. 3 and 5 of Th. 1. If \((\mathrm{M}_{n})\) is a sequence of measurable sets and M is their union, then \(\varphi_{M} = \sup_{n} \varphi_{M_{n}}\) is measurable by Cor. 1 of Th. 2, whence the corollary.

In particular, since the open sets are measurable:

COROLLARY 3. --- The Borel sets in X (GT, IX, §6, No.~3, Def. 4) are \(\mu\) -measurable for every measure \(\mu\) on X.

